% Separate post-CP17 research. Proved after full root/B/C/D actual-file and source audits.
\subsection{Two-generator charts with imperfect residue fields}

The coefficient field in this subsection is denoted by $\kappa$.
It need not be perfect.  The charts below fix a finite residue
extension, and make no claim about the poles of their coefficients
when $\kappa$ is a function field.

\begin{lemma}[Proved: a primitive residue element in transcendence degree one]
\label{D:one-variable-residue-primitive}
Let $k$ be perfect, let $\kappa/k$ be a finitely generated field
extension of transcendence degree one, and let $\lambda/\kappa$
be finite.  Then $\lambda=\kappa(\overline u)$ for some
$\overline u\in\lambda$, even when $\lambda/\kappa$ is inseparable.
\end{lemma}

\begin{proof}
Only characteristic $p>0$ requires discussion.  Write $\lambda_s$
for the maximal separable subextension and
$[\lambda:\lambda_s]=q=p^a$.  All the fields in question have
$p$-rank one.  A purely inseparable extension of degree $q$ is
contained in $\lambda_s^{1/q}$, and the latter also has degree
$q$.  Thus $\lambda=\lambda_s^{1/q}$.
Choose a primitive element $v_0$ of the separable extension
$\lambda_s/\kappa$ and $z\in\kappa\setminus\kappa^p$.
Separability gives $dz\ne0$ in $\Omega_{\lambda_s/k}$.
At least one of $v_0$ and $v_0+z$ has nonzero differential.
Call it $v$.  Translation by $z\in\kappa$ preserves primitivity,
and $dv\ne0$ says $v\notin\lambda_s^p$.
If $a>0$, the element $\overline u=v^{1/q}$ generates an extension
of degree $q$ over $\lambda_s$, and $\kappa(\overline u)$ contains
$\lambda_s=\kappa(v)$.  It is therefore $\lambda$.
For $a=0$ use $v$ itself.
\end{proof}

\begin{proposition}[Proved: finite flat sparse charts and exact finite isomorphism witnesses]
\label{D:sparse-DVR-linear-atlas}
Fix a field $\kappa$, integers $e,f\ge1$, a monic irreducible
polynomial $\psi\in\kappa[U]$ of degree $f$, and an integer
$\delta\ge0$.  Put
\[
 n=ef,\qquad N=\left\lfloor\frac{2\delta}{e}\right\rfloor+2.
                                                        \tag{D.SD.1}
\]
There is an affine polynomial chart over $\kappa$ with at most
$(3N+2)n$ coordinates, and a polynomial witness for equivalence
of two chart objects with at most $(8N+4)n+1$ coordinates
including both objects.  All equations may be taken of degree
at most
\[
                         (4N+12)n^3+1.                 \tag{D.SD.2}
\]
The assertions have the following precise meaning.

For every $\kappa$-algebra $R$, a chart point gives a finite free
$R[[t]]$-algebra $C$ of rank $n$, with elements $u,s$ such that
$s^e=t$ times a unit, $s$ is a nonzerodivisor,
$C/(s)=R[U]/(\psi)$, and $C[1/t]$ is finite etale over
$R((t))$.  Its relative K\"ahler different is $(s^\delta)$.
Every $\kappa$-valued object is a complete DVR with uniformizer
$s$, ramification index $e$, residue field $\kappa[U]/(\psi)$,
and different exponent $\delta$.

For two chart points over any $R$, the finite witness has an
$R$-point if and only if their complete algebras are continuously
isomorphic over $R[[t]]$.  This is an existence test, not an
assertion that a prescribed truncated isomorphism lifts without
loss of precision.  No marking of the residue algebra or
augmentation is imposed on the isomorphisms.

Finally, every finite complete DVR $B$ over $\kappa[[t]]$ with
separable fraction field, these invariants, and residue field
generated by a root of $\psi$, is isomorphic to a chart object.
\end{proposition}

\begin{proof}
\emph{The two equations and their topology.}
Let
\[
 \begin{split}
 P&=U^f-\sum_{0\le i<f,\,0\le j<e} a_{ij}(t)U^iS^j,\\
 Q&=S^e-t\sum_{0\le i<f,\,0\le j<e} b_{ij}(t)U^iS^j,
 \end{split}                                            \tag{D.SD.3}
\]
where all $a_{ij},b_{ij}$ have degree less than $N$ in $t$.
Impose $P(0,U,0)=\psi(U)$.  Write $V=\sum b_{ij}U^iS^j$.
Add $f$ coordinates for an inverse of $V(0,U,0)$ in
$R[U]/(\psi)$, and impose the corresponding $f$ product equations.
Define
\[
 C=\varprojlim_m R[t,U,S]/(t^m,P,Q).
                                                        \tag{D.SD.4}
\]
This is the $t$-adic completion, with $U$ an ordinary polynomial
variable.  In particular, $U$ is not assumed topologically
nilpotent.  Over $\kappa$ the chosen irreducible $\psi$ specifies
the unique residue field; over a general $R$ we retain the whole
finite residue algebra $R[U]/(\psi)$.

We prove, including nonreduced coefficient rings, that
\[
                 \{u^is^j:0\le i<f,\ 0\le j<e\}
                                                        \tag{D.SD.5}
\]
is a basis over $R[[t]]$.  First work over the Noetherian universal
coefficient ring $R_0$ for the displayed finitely many coefficients
and inverse equations.  In $R_0[[t]][U,S]$, the sequence
$(t,P,Q)$ is regular.  Indeed $t$ is regular, $\overline P$
is monic in $U$, its quotient is free over $R_0[S]$, and
$\overline Q=S^e$ is regular on that quotient.
At primes containing $(t,P,Q)$, regular sequences permute
\citep[Lemma~10.68.4, complete proof, Tag 00LJ]{StacksSparseDVRB}.
Consequently multiplication by $t$ on the quotient by $(P,Q)$
has kernel supported away from $V(t)$.  Completion is exact on
finite modules over this Noetherian ring, so the completed
quotient $C_0$ is $t$-torsion free.

Modulo $t$, the box (D.SD.5) is a basis over $R_0$:
first reduce by $S^e$, and then by the monic polynomial
$\overline P$ in $U$.  Successively lifting coefficients modulo
$t,t^2,\ldots$ gives a surjection
$R_0[[t]]^n\to C_0$.  The $t$-torsion freeness shows that its
kernel reduces to zero modulo $t$.  The kernel is finite,
and $t$ lies in the Jacobson radical of $R_0[[t]]$; Nakayama
therefore kills it.  This proves the basis assertion over $R_0$.
Modulo each $t^m$ the same quotient commutes with arbitrary
coefficient base change $R_0\to R$.  Passing to the inverse
limit proves (D.SD.5) for every $R$.  More precisely,
$C=C_0\otimes_{R_0[[t]]}R[[t]]$, since the finite free box basis
and its multiplication constants have these reductions.
Thus this is a genuine all-ring flatness proof, independent
of the later field-point DVR argument.

The element $V$ is a unit modulo $(t,S)$ by construction,
hence modulo $t$ because $S$ is nilpotent there, and hence
in $C$ by completeness.  Thus $s^e=tV$.
Multiplication by $t$ is injective on the free module $C$,
so $s$ is a nonzerodivisor, and the $s$-adic and $t$-adic
topologies are cofinal.  The same identity shows
$C/(s)=R[U]/(\psi)$.

\emph{Jacobian stratum and different.}
Put
\[
                 J=\det\frac{\partial(P,Q)}{\partial(U,S)}.
                                                        \tag{D.SD.6}
\]
Add $nN$ coefficients for an element $W$ of $C/t^NC$, and
$f$ coefficients for an inverse of its reduction in $C/(s)$.
Require
\[
                         J=s^\delta W\pmod {t^NC}.
                                                        \tag{D.SD.7}
\]
Every equation here is an equation in the box basis, using
the reduction procedure below.  Because $eN>2\delta$,
the actual Jacobian in $C$ is $s^\delta$ times a unit.
The equations present the relative differential module by
the two-by-two Jacobian matrix.  One may verify this first
over the Noetherian universal ring: finite differential
modules commute with completion, and the usual conormal
sequence applies before completion.  The assertion over
every $R$ then follows by the finite free base change just proved.
Hence
\[
                  \operatorname{Fitt}_0\Omega_{C/R[[t]]}
                                      =(s^\delta).       \tag{D.SD.8}
\]
After inverting $t$, the Jacobian is invertible, so the
finite free algebra is etale.  Over the universal ring,
the fibres at $t=0$ are the zero-dimensional complete
intersections $(\overline P,S^e)$; the other fibres are etale.
Finite flatness and the fibre criterion therefore make this
a finite syntomic morphism
\citep[Lemma~49.11.1(3), Tag 0FKX]{StacksSparseDVRB}.
This property persists under all the base changes above.
Over a DVR the usual different equals this K\"ahler different,
by the finite syntomic Jacobian formula
\citep[Lemmas~49.12.1--49.12.3 and their proofs,
Tag 0BWB]{StacksSparseDVRB}.

For $R=\kappa$, $C/(s)$ is a field and $s$ lies in the Jacobson
radical.  Thus $C$ is local with maximal ideal $(s)$.
It is finite free over the one-dimensional DVR $\kappa[[t]]$,
so has dimension one.  It is a regular local ring, hence a DVR.
The identity $s^e=tV$ and (D.SD.8) give the asserted invariants.
This assertion about field points is separate from flatness
over arbitrary coefficient rings.

\emph{A Newton estimate over arbitrary coefficient rings.}
We will use the following elementary two-variable form.
In a complete ring in which $s$ is regular, suppose a polynomial
pair $F$ has error in $(s^a)$ and its Jacobian determinant at
the chosen point is $s^\delta$ times a unit, where $a>2\delta$.
The adjugate formula gives a Newton correction in
$(s^{a-\delta})$.  Polynomial Taylor expansion, without
division by factorials, puts the new error in
$(s^{2(a-\delta)})$.
The change in the Jacobian determinant is divisible by
$s^{a-\delta}$, so it remains $s^\delta$ times a unit.
Iteration converges, because
$a\mapsto2(a-\delta)$ increases to infinity.
It produces an exact zero differing from the initial point
by $(s^{a-\delta})$.  Regularity of $s$ justifies division by
$s^\delta$; the proof allows arbitrary nilpotents in $R$.

\emph{Every DVR extension occurs.}
Let $B$ be as in the statement.  Choose a uniformizer $s$ and
a lift $u$ of the specified residue generator.  The box
$u^is^j$ is a basis over $\kappa[[t]]$: its reduction modulo
$t$ is a basis by the filtration by powers of $s$, whose
successive residue quotients have dimension $f$ over $\kappa$.
Expand $u^f$ and $s^e/t$ in that basis.  This gives (D.SD.3)
with power-series coefficients, $P(0,U,0)=\psi$, and $V$ a unit.
The preceding basis proof also works with these coefficients
over $\kappa[[t]]$, and the resulting map to $B$ is an
isomorphism modulo $t$, hence an isomorphism of finite free
modules.  Thus these two equations present $B$ topologically.

For completeness, the finite flat algebra $B$ is syntomic:
its closed fibre has the just proved complete-intersection
presentation, and its generic fibre is separable.
The Jacobian of the two equations consequently has order
$\delta$ by the same different formula.
Truncate all coefficient series modulo $t^N$.
At $(u,s)$ the errors have order at least $eN$, and the
Jacobian still has order $\delta$.
Newton gives exact roots of the truncated equations differing
from $(u,s)$ by $(s^{eN-\delta})$.
Here $eN-\delta\ge e$, so the box basis is unchanged modulo
$t$.  The truncated model maps isomorphically to $B$.
Its Jacobian is $s^\delta$ times a unit, so a choice of $W$
and of the two finite inverse witnesses puts it in our chart.

\emph{Finite witnesses for continuous isomorphism.}
For two chart objects $C_1,C_2$, introduce two box elements
$\eta,\beta\in C_1/t^NC_1$, with $nN$ coefficients each.
Require
\[
                     P_2(\eta,\beta)=Q_2(\eta,\beta)=0
                                             \pmod {t^N}.\tag{D.SD.9}
\]
Form the $n$ by $n$ matrix whose columns are the box expansions
of $\eta^i\beta^j$, in the order (D.SD.5).
Add one inverse variable for the constant coefficient in $t$
of its determinant.  Thus (D.SD.9) defines an isomorphism
\[
                         C_2/t^NC_2\simeq C_1/t^NC_1.    \tag{D.SD.10}
\]
Conversely every such truncated isomorphism has these witnesses.

Lift $\eta,\beta$ to $C_1$ by their box polynomials.
The Fitting ideal of the differential module of (D.SD.10)
is invariant under this isomorphism.  Therefore, writing
$s=s_1$, the ideals generated modulo $t^N$ by
$J_2(\eta,\beta)$ and $s^\delta$ agree.
In $C_1$ write
$J_2(\eta,\beta)=s^\delta a+t^N c$ and
$s^\delta=J_2(\eta,\beta)b+t^N d$.
Since $t^N$ is $s^{eN}$ times a unit, cancellation of the
nonzerodivisor $s^\delta$ shows that the relevant coefficient
$a+s^{eN-\delta}c'$ has a multiplicative inverse modulo
$s^{eN-\delta}$.  It is a unit by completeness.
Thus $J_2(\eta,\beta)$ is actually $s^\delta$ times a unit.
The Newton estimate, with $a=eN$, now corrects the two
approximate images to exact roots.  Their corrections lie
in $(s^{eN-\delta})\subset(t)$.
Their box-basis matrix remains invertible modulo $t$,
so the resulting continuous $R[[t]]$-algebra map is an
isomorphism.  Conversely any continuous isomorphism truncates
to (D.SD.9)--(D.SD.10).
This proof imposes neither a pointed condition on $\beta$
nor a prescribed residue-field automorphism.

\emph{Number and degrees of equations.}
There are $2nN$ coefficient variables in (D.SD.3), $f$ for
the inverse of $V$ modulo $s$, $nN$ for $W$, and $f$ for
its inverse modulo $s$.  This gives $(3N+2)n$ as an upper
bound.  Two objects, the two images, and the determinant
inverse give $(8N+4)n+1$.  All identities amount to
$O(Nn)$ box coefficient equations.

Here is an explicit degree bound, so that finite type alone
is not being used quantitatively.  Reduce a monomial modulo
$t^N$ by using $Q$ whenever its $S$-exponent is at least $e$,
then using $P$ when its $U$-exponent is at least $f$.
Every $Q$ step increases the $t$-exponent, so there are at
most $N-1$ such steps in a surviving term.
Each $P$ step lowers the $U$-exponent by at least one;
a $Q$ step raises it by at most $f-1$.
A term of initial $U$-degree $A$ thus uses at most
$A+(N-1)(f-1)$ steps of type $P$ and at most $N-1$
steps of type $Q$.  Each step adds at most one to its
degree in the coefficient parameters.

Before reduction, the substituted equations and box-matrix
entries have $U$-degree at most $n^2$, and degree in
their image parameters at most $n+1$.
Their reduced coefficients therefore have degree at most
$n^2+(N+1)n+1$.
The $n$ by $n$ determinant and its inverse equation have
degree at most $n^3+(N+1)n^2+n+1$.
The Jacobian, unit, and stratum equations have no larger
degree than the more generous common bound (D.SD.2).
This proves all the stated coordinate and degree estimates.
\end{proof}

\begin{corollary}[Proved: linear local coordinates under a normalized different bound]
\label{D:sparse-DVR-normalized-different-complexity}
Fix $\Lambda\ge0$.  In the preceding proposition, if
$\delta/e\le\Lambda$, the object and isomorphism-witness
spaces have $O_\Lambda(n)$ coordinates and equations of
degree $O_\Lambda(n^3)$.  Consequently the number of their
geometric irreducible components is at most
$n^{O_\Lambda(n)}$ (with $n=1$ handled by a fixed constant).
For a one-variable residue field over perfect constants,
Lemma~\ref{D:one-variable-residue-primitive} supplies the
required primitive residue element for every finite DVR
extension.
\end{corollary}

\begin{proof}
The cutoff in (D.SD.1) is at most $2\Lambda+2$.
Apply the explicit preceding estimates and the same elementary
affine component Bezout bound used in
Theorem~\ref{D:uniform-local-representation-growth}.
The residue lemma does not require its extension to be separable.
\end{proof}

\begin{remark}[Unresolved global input]
These are charts over the fixed residue field $\kappa$, for a
fixed residue polynomial $\psi$.  When $\kappa=k(E)$, none of
the estimates bounds the poles on $E$ of $\psi$ or of the
coefficient functions $a_{ij},b_{ij}$.
Nor does a component estimate for these local witness schemes
bound a constrained global extension locus.
Constructing a residue-cover and coefficient presentation
with controlled global poles, and effective compatibility
with the exceptional neighbourhood, remains necessary.
No finite-generation conclusion is drawn from this subsection.
\end{remark}
